% ============================================================================
\clearpage
\section*{v4.9.0-r37 Source-Only: Triality as Self-Traversal, and the Periods of the Structure}
\addcontentsline{toc}{section}{v4.9.0-r37 Source-Only: Triality as Self-Traversal, and the Periods of the Structure}
\label{sec:v490-r37-triality-self-traversal}
% ============================================================================

\paragraph{Purpose.}
Layer v4.9.0-r36 showed that inside \(\OO\) with \(e_7\) fixed every \(\mathbb Z_3\) symmetry is a colour transformation, so the
generations need a second three-fold structure.  The author states (2026-10-23) that the triole was meant as another
representation of the Sphaera, and that the fractal self-traversal transforms between such representations; the number of steps
is not to be fixed at three.  This layer records the pre-registered test H22 (note
\texttt{llm/Vorab\_Notiz\_H22\_Triality\_Selbstdurchwanderung.md}, commit \texttt{ed74608}, fixed before computing) on the triality
form of v4.8.0-r14 and on the periods available to the structure.  Earlier layers are not rewritten.

\noindent\textbf{Source-only status.}
This v4.9.0-r37 layer is source-only.  It does not build a PDF, does not assign a new DOI,
and does not authorize a public release.  The v4.9.0 Version DOI
\href{https://doi.org/10.5281/zenodo.23534954}{10.5281/zenodo.23534954} and the Concept DOI
\href{https://doi.org/10.5281/zenodo.19210728}{10.5281/zenodo.19210728} are unchanged by it.

\subsection*{v4.9.0-r37 results}
\label{subsec:v490-r37-results}

Let \(t(x,y,z)=\operatorname{Re}((xy)z)\) on \(\OO^3\),
\(\mathfrak g=\{(A,B,C)\in\mathfrak{so}(8)^3: t(Ax,y,z)+t(x,By,z)+t(x,y,Cz)=0\}\), \(\tau(x,y,z)=(y,z,x)\) and
\(\rho(x,y,z)=(\bar x,\bar z,\bar y)\); \(T\), \(R_{\pi/2}\) act on the \(\mathbb C\) factor as in v4.9.0-r30.

\paragraph{Theoremlet v4.9.0-r37.1.}
\begin{enumerate}[leftmargin=2em]
\item[(a)] \(\dim\mathfrak g=28\); each projection \(\mathfrak g\to\mathfrak{so}(8)\) is an isomorphism, and the three representations
of \(\mathfrak g\) on the three slots are irreducible and pairwise inequivalent.  Thus \(t\) is the \(\mathrm{Spin}(8)\) triality form.
\item[(b)] \(\tau\) maps \(\mathfrak g\) to itself and permutes the three slot representations cyclically; it is an outer
automorphism of order three.
\item[(c)] The \(\tau\)-fixed elements \((A,A,A)\in\mathfrak g\) form the derivation algebra of \(\OO\), of dimension \(14\); those
with \(Ae_7=0\) form \(\mathfrak{su}(3)\).
\item[(d)] On \(\OO^4\) the five bracketings of \(\operatorname{Re}(xyzw)\) fall into two classes,
\(\operatorname{Re}(((xy)z)w)=\operatorname{Re}((xy)(zw))=\operatorname{Re}(x(y(zw)))\) and
\(\operatorname{Re}((x(yz))w)=\operatorname{Re}(x((yz)w))\).  The cyclic shift of the four slots exchanges the classes, so no single
bracketing is invariant, while their sum is.
\item[(e)] \(\tau,\rho,T,R_{\pi/2}\) generate on \(\mathbb C\otimes\OO^3\) a group of order \(48\), isomorphic to \(S_3\times D_4\), with
element orders \(1,2,3,4,6,12\); \(\tau\) and \(\rho\) preserve \(t\).
\end{enumerate}

\paragraph{Proof sketch.}
(a)--(c) by solving the linear conditions on \(\mathfrak{so}(8)^3\) (\(512\) equations, \(84\) unknowns) and computing ranks,
intertwiner spaces and commutants (script \texttt{llm/k78\_test\_h22.py}, deviations at most \(10^{-14}\)); \(\tau\) preserves \(t\)
because \(\operatorname{Re}((xy)z)=\operatorname{Re}((yz)x)\) (v4.8.0-r14.2), and an automorphism permuting inequivalent
representations is not inner.  (d) from \(\operatorname{Re}((ab)c)=\operatorname{Re}(a(bc))\) and \(\operatorname{Re}(ab)=\operatorname{Re}(ba)\),
checked numerically; the exchange of the classes was found in a check after the evaluation.  (e) by enumerating the group.
\(\square\)

\paragraph{Verdict and reading.}
H22a, b, c and e are confirmed.  H22d, the statement that no single bracketing is cyclically invariant, is confirmed, but by (d)
it gives no prohibition of a fourth slot.
\begin{enumerate}[leftmargin=2em]
\item The second three-fold structure required by v4.9.0-r36 exists: it is triality.  The three slots of \(\OO^3\) carry the three
inequivalent eight-dimensional representations of \(\mathrm{Spin}(8)\), the shift of the slots is the outer triality automorphism,
and what it fixes is \(G_2\supset SU(3)_c\).  In this reading generation is what the traversal permutes and colour is what it
fixes; they are orbit and fixed group of one motion, no longer one object.  It matches the author's reading of the triole as
three representations of the same \(\OO\cong\mathbb R^8\).  That the slots \emph{are} the physical generations is not shown.
\item That there are exactly three rests on the known classification: \(\mathrm{Spin}(8)\) has exactly three eight-dimensional
irreducible representations and \(\operatorname{Out}(\mathrm{Spin}(8))=S_3\).  (a) is consistent with it.  It excludes a fourth slot
with a new eight-dimensional representation, not a repetition of one of the three.  The structure therefore does not by itself
forbid a fourth generation.
\item The structure has no distinguished period.  Period two belongs to \(\rho\), \(T\), \(J_B\) and the alternating triolic shift of
\S\ref{sec:alternating-triolic-orbit-scan}, period three to triality, period four to the duality rotation \(R_{\pi/2}\), and six and
twelve to their combinations.  In the author's picture, the period seen by a readout is a choice of subgroup.
\item The word triole names at least three structures in the corpus: the frame triole matter--photon--tachyon of the original
paper (spring theorem; \(\mathbb Z_2\) with fixed point, amplitudes \(1:1:3\)), the kernel triole with \(120^\circ\) phases (a colour
rotation inside \(\OO\), v4.9.0-r36), and the slot triole of \(\OO^3\) (triality).  The double uses found in v4.8.0-r13 and
v4.9.0-r34 to r36 arise where these are identified.
\end{enumerate}

\subsection*{v4.9.0-r37 status in the three frames}
\label{subsec:v490-r37-frames}

\begin{center}
\small
\begin{tabular}{@{}>{\raggedright\arraybackslash}p{0.30\textwidth}>{\raggedright\arraybackslash}p{0.62\textwidth}@{}}
\toprule
result & M / T / P\\
\midrule
r37.1(a)--(d) & real octonionic statements without the \(\mathbb C\) factor; the same in M, T and P\\
r37.1(e) & contains the frame maps themselves; \(T\) and \(J_B\) exchange M and T, \(R_{\pi/2}\) is the duality rotation; P is the common fixed point\\
\bottomrule
\end{tabular}
\end{center}

\subsection*{v4.9.0-r37 register entries}
\label{subsec:v490-r37-register}

\begin{center}
\small
\begin{tabular}{@{}>{\raggedright\arraybackslash}p{0.34\textwidth}>{\raggedright\arraybackslash}p{0.58\textwidth}@{}}
\toprule
statement & status after v4.9.0-r37\\
\midrule
triality form \(t\) of v4.8.0-r14 & the \(\mathrm{Spin}(8)\) triality form (r37.1(a), (b))\\
separation of colour and generation & colour in the fixed group \(G_2\supset SU(3)_c\) of triality, generation in its orbit (r37.1(c)); physical identification open\\
exactly three generations & not forbidden by the structure alone; the classification excludes a fourth \emph{distinct} slot only (r37.1(d))\\
periods of the self-traversal & \(1,2,3,4,6,12\) (r37.1(e)); no distinguished period\\
\bottomrule
\end{tabular}
\end{center}

\subsection*{v4.9.0-r37 ledger and guard update}
\label{subsec:v490-r37-ledger-guard}

\[
R_{\rm tr}^{v490\text{-}r37}=(1,0,1)
\]
for (triality carrier of the three slots found; prohibition of a fourth generation not obtained; physical identification of the
slots open).

The layer spends no guard:
\[
R_{\rm glob}^{(v490\text{-}r37,+)}=(0,1,1,1,1),
\qquad
\Gamma_{\rm guard}^{v490\text{-}r37}=(1,1,1,1,1),
\qquad
\beta_4=1/256\ \text{unspent}.
\]
It does not imply exact finite \(C5\) zero closure, exact all-readout terminality,
\(C6/T2\), a global photon-clock/RH theorem, or global all-readout closure.
